\documentclass{article}
\usepackage[T1]{fontenc}
\usepackage{lmodern}
\usepackage{graphicx}
\usepackage{amsmath,amssymb}
\usepackage[a4paper,margin=2cm]{geometry}
\usepackage[absolute,overlay]{textpos}
\setlength{\TPHorizModule}{1mm}
\setlength{\TPVertModule}{1mm}
\usepackage[indonesian]{babel}
\usepackage{hyperref}
\setlength{\emergencystretch}{2em}
% unit: Halaman 22

\begin{document}

% === Halaman 22 (llm) ===

\section*{NoStega: Noiseless Steganography}

\begin{itemize}
    \item Teknik baru steganografi, ditemukan oleh Desoky (2012)
    \item Tidak membutuhkan \textit{cover} untuk menyembunyikan pesan
    \item Latar belakang: penyembunyian pesan di dalam \textit{cover} dapat membuat kualitas \textit{cover} menjadi terdegradasi $\implies$ dapat diserang secara steganalisis untuk menemukan \textit{embedded message}

    \item \textit{NoStega} melakukan kamuflase dengan cara menyembunyikan pesan dalam bentuk \textit{cover} yang terlihat alami sehingga tidak menimbulkan kecurigaan.

    \item \textit{NoStega} menggunakan berbagai materi untuk melakukan kamuflase seperti grafik, email, game, catatan, dan lain-lain.
\end{itemize}

\end{document}
